% GENERATED FILE. Edit canonical lessons, then run npm run generate:books.
% canonical-source-hash: fnv1a64:6d2b9beb5bae223a
% canonical-lessons: ZH-C132-qian, ZH-W132-qian-thousand, ZH-C132-wan, ZH-C132-ling, ZH-C132-yigong, ZH-C132-shu

\chapter{Thousands}
\label{ch:zh-thousands}
\hlchaptermodality{132}

\begin{chapteropening}
\textbf{By the end of this chapter:} \emph{I can count in thousands and ten thousands, say zero, and say how many in all.}

Complete the last lesson of chapter 132: I can count in thousands and ten thousands, say zero, and say how many in all.
\end{chapteropening}

\section[\texorpdfstring{\zh{千} (qiān)}{千 (qiān)}]{\zh{千} (qiān) — thousand}

\label{lesson:ZH-C132-qian}



\begin{quote}
Before the new one: say the Mandarin for winter, then the Mandarin for season.
\end{quote}

\subsection*{You'll want to know: \zh{千}}
\textbf{\zh{千}} — \emph{qiān} — \textquotedblleft{}thousand\textquotedblright{}.

\begin{quote}
\textbf{\zh{我们大学有三千个学生}} \emph{wǒmen dàxué yǒu sānqiān ge xuésheng} — Our university has three thousand students.
\end{quote}

The new character, in pieces: \zh{千} — a short fall over \zh{十}. The next lesson puts it in your hand.

\begin{cousinweb}
\textbf{\zh{千}} (\emph{qiān}) is inherited from Middle Chinese \emph{tshen} (Guangyun, in Baxter's transcription), glossed 'ten hundreds'.
\end{cousinweb}

\begin{grammarlens}[title={What you've built}]
One more word for counting.
\end{grammarlens}

\subsection*{Your turn}
\begin{itemize}
\raggedright
  \item \emph{Say it:} \emph{qiān}
  \item \emph{Say it:} \emph{qiān}, once more
  \item \emph{Say it:} \emph{wǒmen dàxué yǒu sānqiān ge xuésheng} — Our university has three thousand students.
  \item \emph{Recall:} say the Mandarin for winter, then the Mandarin for season, then say \emph{qiān} again
\end{itemize}

\begin{tcolorbox}[breakable,colback=teal!4,colframe=teal!35!black,title={Before you move on}]
What does \zh{千} mean? (\textquotedblleft{}Thousand\textquotedblright{}.) Say it once more.
\end{tcolorbox}
\section[\texorpdfstring{\zh{千} (qiān)}{千 (qiān)}]{\zh{千} — a short fall over \zh{十}}

\label{lesson:ZH-W132-qian-thousand}



\begin{quote}
The word first, then the character inside it.

\begin{itemize}
\raggedright
  \item \emph{Recall:} say \emph{qiān} — \textbf{R1}, one lesson back, before it decays
\end{itemize}

\begin{quote}
\zh{十} \zh{百}
\end{quote}

\emph{Read it:} each one aloud, with nothing said first — \emph{shí}, \emph{bǎi}
\end{quote}

\subsection*{Script — a short fall over \zh{十}}
\hlblockfigure{\detokenize{figures/ZH-W132-qian-thousand-filmstrip.pdf}}{How \zh{千} is written, stroke by stroke}

\begin{quote}
\zh{千}
\end{quote}

\textbf{Two pieces, top to bottom.}

\textbf{A short left-falling stroke} at the top, from the right down to the left.

\textbf{\zh{十}, two strokes}: the horizontal across, then the vertical down through it.

1 + 2 = three.

\subsection*{Your turn}
\begin{itemize}
\raggedright
  \item \emph{Write it:} \zh{千}, all three, with the model in front of you
  \item \emph{Cover:} write it again from memory, and count the pieces, not the strokes
\end{itemize}

\begin{tcolorbox}[breakable,colback=teal!4,colframe=teal!35!black,title={Before you move on}]
How many strokes in \zh{千}? (\textbf{Three}.) Which number you know sits under the fall? (\textbf{\zh{十}}.) Which comes first, the bar or the vertical? (\textbf{The bar}.)

Source: \href{https://raw.githubusercontent.com/chanind/hanzi-writer-data/68d10a4b21150cae5e1ebbd223eed289cf32d90c/data/\%E5\%8D\%83.json}{Hanzi Writer Data \zh{千}, pinned stroke paths and medians}.
\end{tcolorbox}
\section[wàn]{wàn — ten thousand}

\label{lesson:ZH-C132-wan}



\begin{quote}
Before the new one: say the Mandarin for season, then the Mandarin for thousand.
\end{quote}

\subsection*{You'll want to know: wàn}
\emph{wàn} — \textquotedblleft{}ten thousand\textquotedblright{}.

\begin{quote}
\emph{wǒ xiǎng yǒu yíwàn kuài qián} — I wish I had ten thousand yuan.
\end{quote}

Its characters come later; for now it is a word for the ear.

The tone of \zh{一} changes: yī is said yí before the fourth tone of wàn.

\begin{cousinweb}
\emph{wàn} is inherited from Middle Chinese \emph{mjonH} (Guangyun, in Baxter's transcription). The dictionary glosses this short way of writing the character 'ten thousands'; under its longer way of writing it gives the name of a dance and of an insect.
\end{cousinweb}

\begin{grammarlens}[title={What you've built}]
One more word for counting.
\end{grammarlens}

\subsection*{Your turn}
\begin{itemize}
\raggedright
  \item \emph{Say it:} \emph{wàn}
  \item \emph{Say it:} \emph{wàn}, once more
  \item \emph{Say it:} \emph{wǒ xiǎng yǒu yíwàn kuài qián} — I wish I had ten thousand yuan.
  \item \emph{Recall:} say the Mandarin for season, then the Mandarin for thousand, then say \emph{wàn} again
\end{itemize}

\begin{tcolorbox}[breakable,colback=teal!4,colframe=teal!35!black,title={Before you move on}]
What does wàn mean? (\textquotedblleft{}Ten thousand\textquotedblright{}.) Say it once more.
\end{tcolorbox}
\section[líng]{líng — zero}

\label{lesson:ZH-C132-ling}



\begin{quote}
Before the new one: say the Mandarin for thousand, then the Mandarin for ten thousand.
\end{quote}

\subsection*{You'll want to know: líng}
\emph{líng} — \textquotedblleft{}zero\textquotedblright{}.

\begin{quote}
\emph{tā yìbǎi líng yī suì le} — He is a hundred and one.
\end{quote}

Its characters come later; for now it is a word for the ear.

\begin{cousinweb}
\emph{líng} is inherited from Middle Chinese \emph{leng} (Guangyun, in Baxter's transcription), glossed 'to fall' and, after an older book, 'a gentle rain'. The sense zero is not given there; how it came is not settled.
\end{cousinweb}

\begin{grammarlens}[title={What you've built}]
One more word for counting.
\end{grammarlens}

\subsection*{Your turn}
\begin{itemize}
\raggedright
  \item \emph{Say it:} \emph{líng}
  \item \emph{Say it:} \emph{líng}, once more
  \item \emph{Say it:} \emph{tā yìbǎi líng yī suì le} — He is a hundred and one.
  \item \emph{Recall:} say the Mandarin for thousand, then the Mandarin for ten thousand, then say \emph{líng} again
\end{itemize}

\begin{tcolorbox}[breakable,colback=teal!4,colframe=teal!35!black,title={Before you move on}]
What does líng mean? (\textquotedblleft{}Zero\textquotedblright{}.) Say it once more.
\end{tcolorbox}
\section[yígòng]{yígòng — altogether; in all}

\label{lesson:ZH-C132-yigong}



\begin{quote}
Before the new one: say the Mandarin for ten thousand, then the Mandarin for zero.
\end{quote}

\subsection*{You'll want to know: yígòng}
\emph{yígòng} — \textquotedblleft{}altogether; in all\textquotedblright{}.

\begin{quote}
\emph{wǒmen yígòng yǒu wǔ ge rén} — There are five of us in all.
\end{quote}

Its characters come later; for now it is a word for the ear.

The tone of \zh{一} changes: yī is said yí before the fourth tone of gòng.

\begin{cousinweb}
\emph{yígòng} is built from \zh{一} 'one' and \emph{gòng}, both inherited from Middle Chinese, \emph{'jit} and \emph{gjowngH} (Guangyun, in Baxter's transcription), which glosses the second 'together' and 'all'.
\end{cousinweb}

\begin{grammarlens}[title={What you've built}]
One more word for counting.
\end{grammarlens}

\subsection*{Your turn}
\begin{itemize}
\raggedright
  \item \emph{Say it:} \emph{yígòng}
  \item \emph{Say it:} \emph{yígòng}, once more
  \item \emph{Say it:} \emph{wǒmen yígòng yǒu wǔ ge rén} — There are five of us in all.
  \item \emph{Recall:} say the Mandarin for ten thousand, then the Mandarin for zero, then say \emph{yígòng} again
\end{itemize}

\begin{tcolorbox}[breakable,colback=teal!4,colframe=teal!35!black,title={Before you move on}]
What does yígòng mean? (\textquotedblleft{}Altogether; in all\textquotedblright{}.) Say it once more.
\end{tcolorbox}
\section[shǔ]{shǔ — to count}

\label{lesson:ZH-C132-shu}



\begin{quote}
Before the new one: say the Mandarin for zero, then the Mandarin for in all.
\end{quote}

\subsection*{You'll want to know: shǔ}
\emph{shǔ} — \textquotedblleft{}to count\textquotedblright{}.

\begin{quote}
\emph{wǒ shǔ le yígòng yǒu bā ge} — I counted: there are eight in all.
\end{quote}

Its characters come later; for now it is a word for the ear.

\begin{cousinweb}
\emph{shǔ} is inherited from Middle Chinese \emph{srjuX} (Guangyun, in Baxter's transcription), which quotes an older book: 'to reckon'.
\end{cousinweb}

\begin{grammarlens}[title={What you've built}]
One more word for counting.
\end{grammarlens}

\subsection*{Your turn}
\begin{itemize}
\raggedright
  \item \emph{Say it:} \emph{shǔ}
  \item \emph{Say it:} \emph{shǔ}, once more
  \item \emph{Say it:} \emph{wǒ shǔ le yígòng yǒu bā ge} — I counted: there are eight in all.
  \item \emph{Recall:} say the Mandarin for zero, then the Mandarin for in all, then say \emph{shǔ} again
\end{itemize}

\begin{tcolorbox}[breakable,colback=teal!4,colframe=teal!35!black,title={Before you move on}]
What does shǔ mean? (\textquotedblleft{}To count\textquotedblright{}.) Say it once more.
\end{tcolorbox}
